
\section{Try Again}

Editorial article in "Folkebladet" for May 11, 1898.

By its conduct the majority has made the whole church-body system so distasteful to a large part of our people, both within and outside ecclesiastical connection, that for our nationality’s part the Lutheran Church has become altogether unattainable by the way of the church body. In other words, the United Church has put division upon its program instead of gathering. By that way it is impossible to attain peace in the Lutheran Church, even if one were to go on for thousands of years.

And yet peace and unity are very important Christian virtues and goods, necessary in some measure for the Church’s thriving.

It is unavoidable that one must begin again from the beginning, that is to say, with the congregation. If there is to be any talk of a united Lutheran Church among the Norwegians in America, then the very foundation must be laid deep in apostolic spirit.

But if one begins to take this matter up and to look seriously into how things stand with the congregation among us, then one finds quite remarkable and sorrowful things.

Both in the so-called Inner Mission and in the ordering of congregations that has followed from it, the interest of the church body has, in by far the most places, been altogether predominant. Far too often the mission pastor has been degraded into an agent of the church body, who has labored with more thought for the church body and for church politics than for the salvation of souls and the right establishment of the congregation. And not only that: many a young pastor who entered the work with sincere zeal in his soul has fallen into the hands of some local boss and church politician who in reality has steered the pastor’s work into the tracks of party politics. The result has been that congregational organization has become of a church-political kind instead of resting upon Christian and apostolic ground.

And as the beginning in many places is askew, so very often the continuation of the work becomes in accordance with that askew beginning. Both pastors and congregations become so interested and so trained in labor for the church body and in church politics that there remain neither time nor desire nor strength for the main thing, which is the edification of the congregation. And by laboring for the standing of the church body and for ecclesiastical interests, the torpid flesh has, besides, this advantage: it escapes from spiritual work and from serious thoughts of eternity. In the end both pastors and members of the congregation place their Christianity in this, that they are zealous men of the church body and fanatics of party.

Or is there any serious man or woman who does not see with piercing sorrow that many of the noblest and best powers in our Lutheran Church are wasted and blunted precisely for this reason?

What then is to be done? How shall we take hold of the matter?

People cry out to us that we must see to it that a church body of the old sort is formed. This, indeed, is Kildahl’s cry to us. We are on the wrong way, he thinks, because we do not form a church body. We are Catholic, he says, because we are not governed by the majority at the representative annual meeting.

But can in truth a new church body do it better than the old church bodies? Must it not increase the confusion, the division, and the church politics?

Our meaning is this: that the true remedy consists not in laboring for a church body but for the congregation. We understand the New Testament in such a way that it does not command us to have or to maintain any representative church body. We have said that neither do we find it forbidden. Church bodies must therefore be judged according to the experience of church history; they cannot be established beforehand as necessary. It is otherwise with the congregation. We must have the congregation, if we would be Christians. It is therefore best to begin with the congregation. If by God’s Spirit we obtain real congregations with God’s power in them, and with the advance of the Kingdom of God as their calling, then it will grow brighter over the churchly future, which is now darkened by church-political clouds.

With the freedom of the congregation will come the peace of the Church.

Georg Sverdrup